\section{Experiments}
\subsection{Setup and Seed Roles}
The primary setting is CIFAR-10 with PreActResNet-18
\cite{krizhevsky2009learning,he2016identity}, a deterministic 45,000/5,000
training/validation split, and $\epsilon=8/255$. All methods use SGD with
momentum 0.9, weight decay $5\times10^{-4}$, batch size 128, mixed precision,
30 epochs, and a triangular learning-rate schedule. The collapse-prone recipe
uses zero-start FGSM, step $8/255$, and peak learning rate 0.3. The ordinary
FGSM-RS control uses a uniform random start, step $10/255$, and peak learning
rate 0.2. Recovery uses PGD-2 with step $4/255$.

We keep experimental roles explicit. Detector development uses CIFAR-10 seeds
17, 23, and 42. Held-out detector evaluation uses seeds 101, 202, 303, 404,
505, 606, and 707. Intervention ablations and baseline comparisons use the
matched set 17, 23, 42, 101, and 202. Non-event controls are reported
separately from event recall. Seeds 17, 23, and 42 are used for threshold
calibration. They appear in intervention ablations but are not counted toward
held-out detector-evaluation recall. The five-seed intervention table is a
matched component ablation rather than an independent held-out intervention
estimate.

Transfer tests change one factor at a time. CIFAR-10/ResNet-18 uses the
unchanged CIFAR-10 threshold on event seeds 303--707 and three non-event runs
at seeds 17, 23, and 42. CIFAR-100 uses PreActResNet-18 and the same 20-update
windows, five-window median, and three-window minimum history. Saved
development traces at seeds 17, 23, and 42 set a one-time dataset-specific
threshold of 0.6332649834303036 before evaluation of held-out event candidates
101, 202, and 303 and two random-start non-events. A separate CIFAR-10
threat-budget stress test retains random starts and uses $\epsilon=16/255$,
FGSM step $20/255$, and peak learning rate 0.2.

The CIFAR-100 calibration uses six development trajectories: one
collapse-prone and one random-start run for each of seeds 17, 23, and 42. Its
original fitting code scans unique finite scores and applies a lexicographic
event-level rule: maximize detected development events, minimize alarm episodes
on stable development runs, prefer shorter median first-warning lead, and then
prefer the larger threshold. Only one of the three collapse-prone development
runs meets the registered CO definition. The fitted threshold is
0.6332649834303036. This dataset-specific procedure differs from the CIFAR-10
window-level Youden-$J$ calibration.

\subsection{Held-Out Detection and Transfer}
The primary held-out CIFAR-10 evaluation contains seven registered events
spanning epochs 13--24, and none uses a threshold-development seed. The frozen
monitor is active within 60 updates before all 7/7 events; the median lead of
the earliest active warning inside that 60-update window is 40 updates. One
event trajectory also contains an earlier alarm episode. The three ordinary
FGSM-RS non-event runs contain no alarm episode. On the same event traces, a
periodic PGD-10 probe every 100 updates detects 6/7 events with a median delay
of 40 updates after the event and makes a median of 73 PGD calls before
failure. The 100-update probe is a retrospective detector comparator applied
to the saved collapse trajectories; it does not alter training. By contrast,
the controlled FastAdv+ reimplementation probes a fixed 128-image validation
batch every 20 training updates and uses the observed accuracy drop to trigger
a transient PGD training block. The two frequencies therefore describe
different protocols and computational budgets.

Table~\ref{tab:transfer} separates observed events from all non-events. With
the CIFAR-10 threshold unchanged, ResNet-18 yields 4/5 detected events but only
2/5 within 60 updates; this is partial unchanged-threshold architecture
transfer, not complete cross-architecture generalization. On CIFAR-100, two of
three planned collapse-prone trajectories meet the event definition and both
are detected after the one-time dataset-specific calibration. The third
collapse-prone trajectory does not meet the registered CO definition and is
therefore counted with the two random-start controls: among these three
non-events there is one alarm episode.

\begin{table*}[t]
\centering
\caption{Detector audit. ``Near'' denotes an active warning within 60 updates.
The last column is the median near-warning lead in updates.}
\label{tab:transfer}
\footnotesize
\setlength{\tabcolsep}{4pt}
\begin{tabular}{llrrrrrr}
\toprule
Dataset / budget & Architecture & Planned & Events & Detected & Near & Non-event runs / alarm episodes & Median near lead \\
\midrule
CIFAR-10, $8/255$ & PreActResNet-18 & 7 & 7 & 7/7 & 7/7 & 3 / 0 & 40 \\
CIFAR-10, $8/255$ & ResNet-18 & 5 & 5 & 4/5 & 2/5 & 3 / 0 & 50 \\
CIFAR-100, $8/255$ & PreActResNet-18 & 3 & 2 & 2/2 & 2/2 & 3 / 1 & 40 \\
CIFAR-10, $16/255$ RS & PreActResNet-18 & 3 & 3 & 3/3 & 3/3 & -- & 40 \\
\bottomrule
\end{tabular}
\end{table*}

\begin{figure*}[t]
    \centering
    \includegraphics[width=0.94\textwidth]{transfer_and_timing.png}
    \caption{Left: event recall across the primary setting, unchanged-threshold
    architecture transfer, dataset-calibrated transfer, and a threat-budget
stress test, with Wilson 95\% intervals. Right: immediate-switch and
    immediate-replay start PGD-2 with zero intervening FGSM updates; next-epoch escalation
    starts after unprotected CO in two of five matched trajectories.}
    \label{fig:transfer-timing}
\end{figure*}

\subsection{Randomized-FGSM Stress Test}
We first test eight randomized-FGSM runs at $\epsilon=8/255$: three standard
runs at peak learning rate 0.2 and five runs at 0.3. All eight remain robust,
and the detector produces zero alarm episodes. The monitor does not intervene
in this tested standard $8/255$ configuration.

Prior work reports that random starts become less reliable at larger budgets
\cite{andriushchenko2020understanding}. At $\epsilon=16/255$, all three
random-start runs meet the CO definition and terminate at 0\% PGD-10. The
frozen CIFAR-10 detector warns 20, 40, and 60 updates before collapse.
The intervention logs identify these three runs as the immediate-replay
ablation, not the immediate-switch main method. Immediate-replay begins PGD-2 before every
event and finishes at 25.55\%, 26.80\%, and 27.11\% PGD-10. Mean clean and
PGD-10 accuracy are
$71.05\pm0.43\%$ and $26.48\pm0.83\%$. This is a threat-budget stress test,
not cross-dataset evidence.

\subsection{Actionable Lead and Component Controls}
Table~\ref{tab:timing} distinguishes statistical lead from usable lead. All
three responses receive the identical first warning. Immediate-switch and
immediate-replay begin PGD-2 with no intervening FGSM update and act before all
five matched collapses. Waiting for the epoch boundary adds a median of 212 FGSM updates;
two collapses occur before PGD-2 starts.

\begin{table*}[t]
\centering
\caption{Five-seed warning-to-action audit (mean $\pm$ sample std.).
Warn$\rightarrow$start counts intervening FGSM optimizer updates.}
\label{tab:timing}
\scriptsize
\setlength{\tabcolsep}{3pt}
\begin{tabular}{lrrrrr}
\toprule
Response & Before CO & Warn$\rightarrow$start & Clean & PGD-10 & Core min \\
\midrule
Immediate-switch (ours) & 5/5 & 0 updates & $83.06\pm0.65$ & $47.05\pm0.74$ & $18.74\pm0.63$ \\
Immediate-replay (ablation) & 5/5 & 0 updates & $83.14\pm0.62$ & $46.83\pm0.55$ & $18.68\pm0.70$ \\
Next epoch & 3/5 & 212 updates & $83.09\pm0.57$ & $46.72\pm1.05$ & $18.14\pm0.58$ \\
\bottomrule
\end{tabular}
\end{table*}

All component controls use seeds 17, 23, 42, 101, and 202. Immediate-switch
differs from immediate-replay by $+0.22$ PGD-10 points (paired bootstrap
95\% CI $[-0.22,+0.66]$) and $+0.06$ core minutes (CI
$[-0.18,+0.30]$); the clean difference is $-0.07$ points (CI
$[-0.16,0.00]$). Neither comparison supports an immediate-replay advantage. Relative
to waiting for the next epoch, immediate-switch starts PGD-2 157.6 updates earlier on
average (CI $[-252.0,-56.8]$), covers 5/5 rather than 3/5 events before CO,
with a proportion difference of $+0.40$ (CI $[0.00,+0.80]$), and differs by
$+0.33$ PGD-10 points (CI $[+0.08,+0.66]$). Its clean difference is $-0.02$
points (CI $[-0.14,+0.08]$), while core time increases by 0.59 minutes (CI
$[+0.47,+0.71]$). The before-CO point estimate favors immediate-switch, but five
seeds are insufficient to establish a clear proportional advantage. The small
endpoint difference is secondary to the intervention-timing result. Full PGD-2
differs from the next-epoch policy by $+0.09$ points (CI
$[-0.27,+0.42]$) but adds 3.46 core minutes. The detector supplies the
warning, PGD-2 supplies recovery, and immediate-switch converts warning
into pre-collapse action; rollback itself is not the active ingredient.

\begin{table}[t]
\centering
\caption{Matched five-seed controls (mean $\pm$ sample std.)}
\label{tab:components}
\scriptsize
\setlength{\tabcolsep}{2.5pt}
\begin{tabular}{lrrr}
\toprule
Policy & Clean & PGD-10 & Core min \\
\midrule
Immediate-switch (ours) & $83.06\pm0.65$ & $47.05\pm0.74$ & $18.74\pm0.63$ \\
Immediate-replay (ablation) & $83.14\pm0.62$ & $46.83\pm0.55$ & $18.68\pm0.70$ \\
Next epoch & $83.09\pm0.57$ & $46.72\pm1.05$ & $18.14\pm0.58$ \\
Fixed epoch 13 & $83.32\pm0.70$ & $46.41\pm0.50$ & $19.37\pm0.21$ \\
Fixed epoch 21 & $82.35\pm1.12$ & $45.75\pm0.94$ & $17.85\pm0.21$ \\
PGD-2 from epoch 1 & $83.57\pm1.02$ & $46.81\pm0.90$ & $21.60\pm0.13$ \\
\bottomrule
\end{tabular}
\end{table}

\subsection{Controlled CO-Mitigation Baselines}
FastAdv+ and GradAlign are controlled reimplementations rather than claims of
using official released code. They share the split, architecture, five seeds,
30-epoch schedule, checkpoint rule, and PGD-10 evaluator. FastAdv+ uses one
fixed 128-image validation batch, a PGD-10 probe every 20 optimizer updates
with step $2/255$, a 0.10 drop from the best observed probe accuracy as its
trigger, and a transient 20-update PGD-10 training block before returning to
random-start FGSM. GradAlign uses $\lambda=0.2$, a uniform perturbation in
$[-\epsilon,\epsilon]$ clipped to the image domain, cosine alignment between
clean and random-point input gradients, and a second-order graph only for the
random branch. It otherwise uses batch size 128 and the shared SGD, momentum,
weight decay, and triangular learning-rate schedule.

Table~\ref{tab:baselines} reports measured core parameter-update time.
FastAdv+'s separately logged online PGD probes add
$3.14\pm0.00$ minutes, giving $30.41\pm3.53$ measured minutes including
those probes. Immediate-replay's snapshot save-restore operations were not
timed, so 18.68 minutes is not a complete end-to-end measurement. GradAlign
achieves the highest mean PGD-10 accuracy among these 30-epoch methods; the
proposed policy occupies a lower-robustness, lower-measured-core-cost point.

\begin{table}[t]
\centering
\caption{Unified five-seed CO-mitigation baselines}
\label{tab:baselines}
\scriptsize
\setlength{\tabcolsep}{3pt}
\begin{tabular}{lrrr}
\toprule
Method & Clean & PGD-10 & Core min \\
\midrule
Ours, next epoch & $83.09\pm0.57$ & $46.72\pm1.05$ & $18.14\pm0.58$ \\
Ours, immediate-switch & $83.06\pm0.65$ & $47.05\pm0.74$ & $18.74\pm0.63$ \\
FastAdv+ & $81.96\pm0.60$ & $47.09\pm0.71$ & $27.27\pm3.53$ \\
GradAlign & $79.54\pm0.86$ & $48.31\pm0.81$ & $45.71\pm0.11$ \\
\bottomrule
\end{tabular}
\end{table}

\subsection{Strong Test-Set Evaluation}
Table~\ref{tab:test} reports full-test clean and PGD evaluation on all 10,000
CIFAR-10 test images. Table~\ref{tab:aa} gives a supplementary
AutoAttack-standard evaluation on
the first 1,000 examples in the unshuffled official test-set order. AA-1k uses
the $\ell_\infty$ threat model at $\epsilon=8/255$ and the standard attack
suite comprising APGD-CE, APGD-T, FAB-T, and Square. The immediate-switch, immediate-replay,
next-epoch, FastAdv+, and GradAlign checkpoints use seed 101; the FGSM-RS and
PGD-10 training references use seed 17. In each run,
\emph{best\_robust.pt} is selected by the highest recorded validation PGD-10
accuracy, before test evaluation. This is a representative checkpoint
comparison rather than a cross-seed aggregate. Validation PGD-10 selects
checkpoints but does not trigger the online monitor. The immediate-switch
seed-101 checkpoint was fixed before
any result in Tables~\ref{tab:test}--\ref{tab:aa} was observed; no alternative
seed or checkpoint was tested for replacement.
The $46.93\%$ PGD-10 entry in Table~\ref{tab:test} is from this single
pre-fixed seed-101 checkpoint. Table~\ref{tab:timing}'s $47.05\pm0.74\%$ is
instead the mean and sample standard deviation across five intervention seeds;
the two summaries need not match.

For the primary immediate-switch method, PGD-50 and five-restart PGD-20 give 45.52\%
and 45.66\%. Their agreement, and the analogous agreement for immediate-replay and the
training references, does not reveal an obvious weak-iterative-attack failure.
On the same fixed 1,000 images, the full standard AutoAttack suite gives
42.30\% for immediate-switch, 42.60\% for immediate-replay, 41.30\%
for next-epoch escalation, 43.00\%
for ordinary FGSM-RS, 48.70\% for PGD-10 training, 42.40\% for FastAdv+, and
44.40\% for GradAlign. The two controlled baselines use their seed-101
checkpoints selected by the same validation rule. AA-1k is a subset estimate,
so it complements the full-test PGD results without ruling out gradient masking.

\begin{table}[t]
\centering
\caption{Full-test strong CIFAR-10 accuracy (\%)}
\label{tab:test}
\scriptsize
\setlength{\tabcolsep}{2pt}
\begin{tabular}{lrrrr}
\toprule
Method & Clean & PGD-10 & PGD-50 & PGD$\times5$ \\
\midrule
Immediate-switch (ours) & 82.95 & 46.93 & 45.52 & 45.66 \\
Immediate-replay (ablation) & 82.97 & 47.20 & 45.46 & 45.65 \\
Next epoch & 82.57 & 46.38 & 44.75 & 44.91 \\
Fast FGSM-RS & 82.48 & 47.77 & 46.24 & 46.36 \\
PGD-10 training & 82.47 & 51.76 & 50.77 & 50.84 \\
\bottomrule
\end{tabular}
\end{table}

\begin{table}[t]
\centering
\caption{AutoAttack-standard accuracy on the fixed first 1,000 CIFAR-10 test
images (\%), $\ell_\infty$, $\epsilon=8/255$}
\label{tab:aa}
\scriptsize
\setlength{\tabcolsep}{6pt}
\begin{tabular}{lr}
\toprule
Method & AA-1k \\
\midrule
Immediate-switch (ours) & 42.30 \\
Immediate-replay (ablation) & 42.60 \\
Next epoch & 41.30 \\
Fast FGSM-RS & 43.00 \\
PGD-10 training & 48.70 \\
FastAdv+ & 42.40 \\
GradAlign & 44.40 \\
\bottomrule
\end{tabular}
\end{table}
